% Приложение: Список литературы (формат ГОСТ)
\chapter{Список литературы}
\label{app:references}

% Основной список оформлен по ГОСТ в виде thebibliography.
% По мере написания глав пополняйте; нумерация сквозная.
\begin{thebibliography}{99}

\bibitem{tannenbaum}
Таненбаум, Э.~С. Архитектура компьютера : [пер. с англ.] / Э.~Таненбаум, Т.~Остин. --- 6-е изд. --- Санкт-Петербург [и др.] : ПИТЕР, 2014. --- 811 с. : ил. ; 24 см. --- (Классика Computer Science). --- ISBN 978-5-496-00337-7.

\bibitem{krahotkina}
Архитектура ЭВМ : учебное пособие / авт.-сост. Е.~В. Крахоткина, В.~И. Терехин ; Северо-Кавказский федеральный университет. --- Ставрополь : СКФУ, 2015. --- 80 с. --- Режим доступа: по подписке. --- URL: \url{https://biblioclub.ru/index.php?page=book\&id=457862}. --- Текст : электронный.

\bibitem{dikov}
Диков, А.~В. Компьютер изнутри : учебное пособие : [16+] / А.~В. Диков. --- Москва ; Берлин : Директ-Медиа, 2015. --- 127 с. : ил., схем. --- ISBN 978-5-4475-5530-6. --- Текст : электронный.

% TODO: пополнять по мере написания глав (shell, Docker, Git, C/компиляция).
\end{thebibliography}

\section{Интернет-ресурсы}
\begin{description}
  \item{Страница курса} --- учебные материалы и задания на Google Classroom
  (\url{https://classroom.google.com}).
  \item{GNU Bash Manual} --- \url{https://www.gnu.org/software/bash/manual/}
  \item{Coreutils documentation} --- \url{https://www.gnu.org/software/coreutils/}
  \item{Docker documentation} --- \url{https://docs.docker.com/}
  \item{Pro Git (online)} --- \url{https://git-scm.com/book}
  % TODO: пополнять.
\end{description}
